\section{Data, measures and research design}\label{sec:data}

\subsection{Why these data, and the sample}\label{sec:sample}

The hypotheses need four things from the data, and the Korean record supplies each in a form that few other settings do. First, the filing that carries the controller's statements must be mandatory, public and comparable across firms and years: every Korean company that divides must file a division resolution with the Financial Supervisory Service's disclosure system, DART, and since 2017 the filing follows a structured form whose items---purpose, effects, the new company's assets and debt, listing application, put-option contracts---are the same for every filer, so the text can be read as a document and its stated financials used as the baseline that \eqref{eq:posterior} requires. Second, the outcome must be a discrete event the parent must later report: a Korean parent lists its equity investments, with the stake held, in every annual and half-year report, and the exchange records every listing, so whether the subsidiary was kept, sold or listed can be read for every division old enough to have reports. Third, the setting must offer an untreated comparison for the reform: equity spin-offs were filed under the same form in the same years and were exempt from all three measures. Fourth, an independent hand-coded sample should exist against which a model-read measure can be checked: \citeA{kang2025} coded the purposes and outcomes of \KLN{} listed-parent divisions of 2001--2021 by hand, and Section~\ref{sec:validity} compares our extractions with theirs. Prices come from the Korea Exchange for 2017--2026 and from Compustat Global for 2001--2016, the latter also supplying the universe of Korean common stocks from which the size and book-to-market benchmarks and the propensity-matched controls are drawn; ownership comes from the largest-shareholder table of the annual report; foreign holdings from the exchange's daily shareholding records.

The population is every \captive{} announced in a division resolution filed between 2001 and 2026, whether or not the parent was listed. For 2017 onward we take the structured division-resolution reports that DART returns (\NNewN{} divisions, \NMixed{} of which combined a captive with an equity spin-off). For 2001 to 2016, before the structured record exists, we retrieved every filing titled a division or division-and-merger resolution, clustered a filing and its corrections into one event, and had a language model extract the division's form with verbatim evidence; \NOldN{} events have a captive part. Divisions are dated by the receipt date of the original filing. The events pass through three measurements (Table~\ref{tab:sample}, Panel A): \NScored{} of the \NAllN{} events have a reading of the filing by both models, \NClean{} a clean outcome label (\NGateOne{} with the filing financials the baseline model needs), \NAllPriced{} listed parents have prices around the announcement and \NScoreOneY{} a full one-year window and a reading, a window being full when at least 90 percent of trading days $+2$ to $+250$ are priced. Panel B describes those \DescN{} parents, announced between \DescYearMin{} and \DescYearMax{} (\DescPeriodOld{} before 2017), \DescKosdaqPct{} percent of them KOSDAQ companies. The signal has a mean of \DescPLlmMean{} and a standard deviation of \DescPLlmSD{}; \DescDisposedPct{} percent of the labelled subsidiaries were disposed of and \DescEarlyRecPct{} percent of those with a recorded timing within a year; the mean three-day announcement return is \DescCarMmMPMean{} percent and the median one-year market-adjusted return \DescBharMadjPPMedian{} percent, against a mean of \DescBharMadjPPMean{}, the skew that motivates the rank specification below; the controlling block averages \DescBlockMean{} percent of common shares and foreign ownership \DescForeignMean{} percent with a median of \DescForeignMedian{}.

\begin{table}[htbp]
\centering
\caption{Sample construction and descriptive statistics}\label{tab:sample}
{\footnotesize\singlespacing\setlength{\tabcolsep}{3pt}
\emph{Panel A. From division filings to the analysis samples}\par\smallskip
% generated by code/make_tables_cgir.py (labels remapped by make_tables_jifma.py) on 2026-11-11 — source: carveout_disposal_tracking_{2001_2016,v4}.json; dart_division_decisions_2017_2026.json; event_study_meta_compustat.csv; ownership_pre_event_all.json; gate1_pooled_20261019.json by_source
\begin{tabular}{>{\raggedright\arraybackslash}p{9.2cm}rrr}
\toprule
 & 2001--2016 & 2017--2026 & Total \\
\midrule
Captive spin-off announcements with a tracked outcome file (listed and unlisted parents) & 464 & 294 & 758 \\
\quad disposed within the tracking window & 257 & 145 & 402 \\
\quad retained at 100\% & 89 & 90 & 179 \\
\quad merged back or liquidated (excluded) & 43 & 29 & 72 \\
\quad label unresolved or too recent (excluded) & 75 & 30 & 105 \\
Clean-label events with a disposal score (Table~\ref{tab:gate1}) & 345 & 234 & 579 \\
Listed parents with daily prices around the announcement & 380 & 260 & 640 \\
\quad with a full one-year post-announcement window & 366 & 249 & 615 \\
Parents with pre-announcement controlling-block ownership & 387 & 283 & 670 \\
\bottomrule
\end{tabular}
\par\medskip
\emph{Panel B. Listed parents with a signal and a full one-year return window}\par\smallskip
{\scriptsize% generated by code/make_tables_jifma.py on 2026-11-14 — source: data/processed/bhar_controls_events_20261025.csv (19), ownership_pre_event_all.json, event_governance_controls_20261110.csv (34); sample = signal and full one-year window
\begin{tabular}{>{\raggedright\arraybackslash}p{6.6cm}rrrrrr}
\toprule
Variable & $n$ & Mean & SD & P25 & Median & P75 \\
\midrule
Disposal signal $S$ (0--1) & 559 & 0.525 & 0.188 & 0.403 & 0.490 & 0.631 \\
Disposed (clean label, 0/1) & 469 & 0.70 & 0.46 & 0.00 & 1.00 & 1.00 \\
Disposed within a year (timing recorded, 0/1) & 405 & 0.48 & 0.50 & 0.00 & 0.00 & 1.00 \\
Announcement CAR $[-1,+1]$, \% & 546 & 0.6 & 8.4 & $-$3.3 & 0.2 & 3.6 \\
One-year market-adjusted BHAR, \% & 559 & $-$5.6 & 70.9 & $-$41.0 & $-$16.1 & 11.5 \\
Within-year rank of the BHAR (0--1) & 559 & 0.52 & 0.29 & 0.27 & 0.52 & 0.77 \\
Controlling block, \% of common shares & 524 & 35.0 & 19.1 & 21.5 & 34.8 & 46.6 \\
Foreign ownership at day $-1$, \% & 540 & 8.0 & 12.9 & 0.4 & 2.4 & 10.0 \\
Log market capitalisation (KRW) & 557 & 25.89 & 1.74 & 24.78 & 25.53 & 26.81 \\
Log book-to-market & 408 & $-$0.40 & 1.00 & $-$0.98 & $-$0.34 & 0.29 \\
Six-month pre-announcement market-adjusted return & 546 & 0.06 & 0.56 & $-$0.22 & $-$0.04 & 0.22 \\
KOSDAQ parent (0/1) & 559 & 0.47 & 0.50 & 0.00 & 0.00 & 1.00 \\
New company's share of combined assets & 317 & 0.13 & 0.12 & 0.04 & 0.10 & 0.19 \\
Surviving company's debt to assets & 324 & 0.41 & 0.21 & 0.25 & 0.41 & 0.57 \\
Listing application stated in the filing (0/1) & 559 & 0.002 & 0.042 & 0.000 & 0.000 & 0.000 \\
\bottomrule
\end{tabular}
}}
\tabnote{Panel A counts events at each stage of measurement, by filing era. Panel B describes the \DescN{} parents of the return regressions; variables are defined in Table~\ref{tab:defs} and Section~\ref{sec:returns}, returns in percent.}
\end{table}

The outcome is read from the parent's own later filings. We traced each new company by name (and by the name it took at incorporation) through the parent's annual and half-year reports for up to six fiscal years, to fiscal 2025 at the latest, and through the exchange's listing records. The subsidiary is disposed of when the parent sold a majority, listed it, or let its stake fall below 99.5 percent through a sale or a third-party investment, and retained when every report through the latest one, fiscal 2025 at the most recent, shows 100 percent; divisions after June 2025, which have no annual report to read, and divisions merged back, liquidated or too poorly documented to classify are excluded (\NAllOther{} and \NAllUnresolved{} events). A retained label is therefore right-censored---a 2024 division has one or two reports behind it, a 2010 division six---and the forecast's horizon of three to five years is met only for divisions up to about 2021; the hazard model of Section~\ref{sec:gate1} treats the censoring explicitly, and the period columns of Table~\ref{tab:gate1} and Table S6 give the older divisions alone. Under the 99.5 percent threshold a small third-party stake counts as a disposal; the labels that turn on it are flipped in Table S2. Of the \NClean{} clean-label events, \CleanDisposedShare{} percent were disposed of and \IPOShare{} percent listed; on listed parents of 2001--2021 our labels place \OurListedShare{} percent of subsidiaries in a listing and \OurRetainedShare{} percent retained at 100 percent, against \KLListed{} and \KLRetained{} percent in the independently built hand-coded sample of \citeA{kang2025}. A disposal is dated by the fiscal year of the report that first shows it, which separates disposals within a year of the division from later ones and gives the survival analysis its time scale. The 2017--2026 labels agree with a hand-maintained tracker on \LabelCrossAgree{} decided cases (\LabelCrossPct{} percent); \NSensitive{} depend on a definitional choice or a source conflict, and recoding all of them as retained changes no result (Table S2).

\subsection{The disposal signal}\label{sec:signal}

We estimate the reader's forecast $\pi_i$ of \eqref{eq:posterior} with two large language models from different vendors, Model~A, Anthropic's \texttt{claude-fable-5-1}, and Model~B, OpenAI's \texttt{gpt-6.1-sol}, both called through Amazon Bedrock under the identifiers \texttt{us.anthropic.claude-fable-5-1} and \texttt{us.openai.gpt-6.1-sol} between 15 and 25 October 2026, with extended reasoning and no sampling parameters altered. Each model reads the filing's full text and is told the filing's own date, so that its forecast is made as of that date; the reading stands for what a careful reader could have inferred then, a proxy whose validity Section~\ref{sec:validity} tests. Text measures of disclosure content have moved from word lists \cite{loughran2011} to supervised learning \cite{frankel2022} and to language models \cite{siano2025}; ours asks for a forecast of a specific event that we later observe, not a tone score. Each model is instructed to forecast, from the document alone and without using anything it may know about the company's later history, the probability that the parent will not retain 100 percent of the new company within roughly three to five fiscal years, and to return it as an integer from 0 to 100 together with the most likely route, reasons, verbatim supporting sentences, and the facts stated in the filing (purposes from a fixed list of seven categories; whether a listing, an outside investor, a sale or buyer, protection measures and the appraisal right are mentioned; the new company's business and capital). The disposal signal of filing $i$ is the two models' probabilities averaged and rescaled to the unit interval,
\begin{equation}\label{eq:signal}
S_i = \frac{1}{2}\sum_{m\in\{A,B\}}\frac{1}{R_m}\sum_{r=1}^{R_m}\frac{p_{i,m,r}}{100},
\end{equation}
where $R_m$ is the number of runs (two for 2017--2026, one for 2001--2016). It is judged against the realised label $D_i$ by the area under the ROC curve, $\mathrm{AUC} = \Pr(S_i > S_j \mid D_i = 1, D_j = 0)$, and enters the return tests as a level or as its within-year rank. The prompts, parsing code and every judgment are in the replication archive, so the tables regenerate without a model call.

\subsection{Measurement validity and bias of the signal}\label{sec:validity}

Because $S$ is produced by a model rather than by a hand-coded rule we treat it as a measured variable with error and ask of it what one would ask of any new measure: whether it is stable, whether it could have been contaminated, for whom it works less well, whether its errors line up with firm characteristics, and whether its ingredients agree with hand coding (Table S1; predictive validity is Section~\ref{sec:gate1}).

Re-running each model reproduces $S$ with correlations of \RepeatCorrA{} and \RepeatCorrB{}. Across vendors the probabilities correlate at \TwoModelCorr{} over the \TwoModelN{} filings with a first-run judgment from both models (three of the \NScored{} scored filings have Model~A's reading from a repeat run only), and the models differ in level (means \TwoModelMeanA{} and \TwoModelMeanB{}), which is why we average and rank within year.

A model that recognises the firm could ``predict'' an outcome it has read about \cite{glasserman2024}, so every filing was re-read with the parent's and new company's names, group names, tickers and dates removed. The anonymised reading has an AUC of \ValAUCAn{} against \ValAUCId{} for the identified first-run reading of the same filings (the two-run average of Section~\ref{sec:signal} reaches \AUCPooledLLM{}); the paired-bootstrap difference is \ValDAUC{} with interval [\ValDAUCLo, \ValDAUCHi], and at most \ValDAUCMaxSub{} in any subgroup. Model~A reported drawing on outside knowledge for \ValSubsetN{} of the \ValOutsideN{} labelled filings for which it returned the field (\ValOutsideA{} percent; Model~B reported it for \ValOutsideB{} percent), and inside that subset the identified and anonymised readings are equally accurate (\ValSubsetAUCId{} against \ValSubsetAUCAn{}). The anonymised signal reproduces the return results (Supporting Information). Anonymisation removes the recall of a firm; it does not remove what a model trained to 2026 knows about how division filings tend to end. A careful human reader in 2010 brought the same kind of knowledge, in smaller measure, and the only test that separates the two is a blind human reading of a sample of filings, which we have prepared but not run (Section~\ref{sec:conclusion}).

Discrimination is similar across exchanges (\ValAUCKosdaqKosdaq{} and \ValAUCKosdaqKospiOther{}) and weaker for filings before 2009 (\ValAUCPeriodEarly{} against \ValAUCPeriodLate{} after 2016), for unlisted parents (\ValAUCListedUnlisted{} against \ValAUCListedListed{}; interaction $p$ = \ValLRDiscListed{}) and for the smallest third of parents by filing assets (\ValAUCSizeSmall{} against \ValAUCSizeMid{}--\ValAUCSizeLarge{}; $p$ = \ValLRDiscSize{}). The signal is conservative everywhere (mean \ValMeanS{} against a realised disposal rate of \ValMeanY{}), most so for small parents. The part of $S$ that the models' own extracted facts, text length and period do not explain (stage-one $R^2$ \ValStageOneRTwo{}) has no standardised association with size, exchange, listing status, relative size or foreign ownership above \ValMaxBeta{}; the one pattern is that the error $S - D$ rises with parent size (\ValErrSizeBeta{}, $p$ = \ValErrSizeP{}). The signal is not inflated for large, listed or foreign-held firms; if anything it under-reads small parents, whose disposal rate is highest, and the return tests carry size, exchange, listing and industry controls.

The two models agree on the facts that enter the signal's decomposition (sale purpose $\kappa$ = \KappaSale{}, fund-raising \KappaFund{}, holding-company restructuring \KappaHolding{}, a named sale or buyer \KappaBuyer{}) and disagree only where a field's boundary is ambiguous (whether the form's ``re-listing application: no'' counts as a mention of listing, $\kappa$ = \KappaListing{}), fields we do not use. Against the hand coding of \citeA{kang2025} on the same population, our listed-parent filings of 2001--2021 show a stated sale purpose in \OurSaleStrict{} percent (\OurSaleBroad{} counting fund-raising) against their \KLSale{}; their near-universal specialisation code (\KLSpecial{} percent) and rare efficiency code (\KLEff{}) are not comparable with a multi-label extraction that records both for almost every filing.

\subsection{Returns, ownership and controls}\label{sec:returns}

For listed parents day 0 is the first trading day on or after the receipt date of the original filing. The announcement return is the market-model cumulative abnormal return over days $-1$ to $+1$, estimated over days $-250$ to $-31$ against the parent's market index; the one-year return is the parent's buy-and-hold return minus the benchmark's over days $+2$ to $+250$,
\begin{equation}\label{eq:bhar}
\mathrm{BHAR}_i = \prod_{\tau=2}^{250}(1+R_{i\tau}) - \prod_{\tau=2}^{250}(1+R_{B\tau}),
\end{equation}
with $B$ the market index in the main specification and the parent's size and book-to-market portfolio or a matched firm in the benchmark tests. Windows with fewer than 90 percent of trading days are dropped. Daily adjusted prices for 2017--2026 come from the Korea Exchange and for 2001--2016 from Compustat Global, whose returns match the exchange's on the overlap (median daily correlation \KRXCompustatCorr{}); \GlitchExcluded{} parents whose Compustat window contains a return beyond the exchange's price limit without a delisting are excluded. Because long-run returns are skewed \cite{barber1997, lyon1999, fama1998, mitchell2000}---Panel B of Table~\ref{tab:sample} puts the mean one-year return at \DescBharMadjPPMean{} percent and the median at \DescBharMadjPPMedian{}---the headline dependent variable is the parent's percentile rank of the one-year return within its announcement year, with means and medians beside it.

The controlling block at the fiscal year-end before the announcement year is $\mathrm{Block}_i = 100\,(H_i - T_i)/N_i$, where $H_i$ is the common shares held by the largest shareholder and the related parties listed with it, $T_i$ the treasury shares within that group and $N_i$ the shares outstanding; for announcements before 2016 two models read the table from the annual report and agree within one percentage point in \OwnWithinOne{} of \OwnPairs{} cases. The median block is \BlockMedian{} percent. The controls are the parent's log market capitalisation at the last month-end before the announcement and its log book-to-market (Compustat Global), its market-adjusted return over the six months before the announcement, its market, the new company's assets relative to the surviving company's, the surviving company's debt-to-assets ratio and whether the filing states a listing application; we add foreign ownership at day $-1$ from the exchange's shareholding records (\ForeignCoverage{} events; median \ForeignMedian{} percent) and, for fiscal years from 2014, the share of outside directors on the parent's registered board (\BoardN{} events; median \BoardMedian{}). Missing controls are set to the sample median with an indicator, and continuous controls are winsorised at the first and ninety-ninth percentiles. The regressions take the form
\begin{equation}\label{eq:main}
y_i = \alpha_{t(i)} + \alpha_{j(i)} + \beta\, S_i + \boldsymbol{\gamma}'\mathbf{x}_i + \varepsilon_i,
\end{equation}
with announcement-year effects $\alpha_t$, GICS-sector effects $\alpha_j$ where stated, and standard errors clustered by parent (two-way by parent and year where stated) beside the HC1 errors of the baseline; $\beta$ estimates $b_A$ of \eqref{eq:market} when $y_i$ is the announcement return and $b_Y$ when it is the one-year return or its rank. Table~\ref{tab:defs} collects the ten measures in one-sentence form; Table S9 in the Supporting Information gives the same ten rows, in the same order, with the full operational rules, sources and exclusions.

\begin{table}[htbp]
\centering
\caption{Definitions of the measures used in this paper (short form)}\label{tab:defs}
{\scriptsize\singlespacing\setlength{\tabcolsep}{3pt}% generated by code/make_tables_cgir.py (labels remapped by make_tables_jifma.py) on 2026-11-14 — source: DEFS list in code/make_tables_cgir.py (one row list; full form S_defs_full)
\begin{tabular}{>{\raggedright\arraybackslash}p{0.6cm}>{\raggedright\arraybackslash}p{3.1cm}>{\raggedright\arraybackslash}p{7.6cm}>{\raggedright\arraybackslash}p{3.4cm}}
\toprule
No. & Term & Definition & Applied in \\
\midrule
D1 & Captive spin-off (\emph{mulchŏk bunhal}) & A division in which the parent receives all shares of the new company; the parent's shareholders receive none. & Table~\ref{tab:sample}; all tests \\ \addlinespace[1pt]
D2 & Disposal & The parent ceases to hold about 100\% of the new company within the tracking window: sale of a majority, listing, or a stake sale or third-party investment taking the stake below 99.5\%. & Tables~\ref{tab:sample}, \ref{tab:gate1}, \ref{tab:returns} \\ \addlinespace[1pt]
D3 & Early disposal & First ownership change recorded in the fiscal year of the division or the next. & Tables~\ref{tab:returns}, \ref{tab:block} \\ \addlinespace[1pt]
D4 & Disposal signal & The probability of disposal that two language models assign to the announcement filing, averaged. & Tables~\ref{tab:gate1}, \ref{tab:returns} \\ \addlinespace[1pt]
D5 & Anonymised condition & The same filing with parent and new-company names, group prefixes and dates masked, and the filing date withheld. & Table~\ref{tab:gate1} \\ \addlinespace[1pt]
D6 & Controlling block & Common shares held by the largest shareholder and related parties, as a percentage of common shares, at the fiscal year-end before the announcement year; treasury shares excluded. & Table~\ref{tab:block} \\ \addlinespace[1pt]
D7 & Announcement CAR & Market-model cumulative abnormal return over trading days $-1$ to $+1$ around the filing date. & Tables~\ref{tab:returns}, \ref{tab:reform} \\ \addlinespace[1pt]
D8 & One-year BHAR & Buy-and-hold return of the parent minus that of its market index over trading days $+2$ to $+250$. & Tables~\ref{tab:returns}--\ref{tab:block} \\ \addlinespace[1pt]
D9 & Within-year rank & Percentile rank of the one-year BHAR among the year's events (0--1). & Tables~\ref{tab:returns}, \ref{tab:block} \\ \addlinespace[1pt]
D10 & Post-reform period & Filings received on or after 1 October 2022. & Table~\ref{tab:reform} \\ \addlinespace[1pt]
\bottomrule\end{tabular}
}
\tabnote{Short form; Table S9 gives the same ten rows, in the same order, with data fields, timing rules and exclusions. Both tables are generated from one row list in the replication archive.}
\end{table}

\subsection{Research design}\label{sec:design}

Several alternative explanations stand between the associations and the readings we give them; Table~\ref{tab:design} pairs each with the device that addresses it, and three of the devices are described here because the tables do not show them.

\emph{Benchmarks and selection.} From Compustat Global we build the universe of Korean common stocks and, for every event, a $5\times5$ size and book-to-market portfolio, a $3\times3$ portfolio, a matched firm nearest in size and book-to-market, and a propensity-matched control drawn from a firm-year logit of the division decision on lagged size, Tobin's $q$, leverage, investment, cash flow, cash, sales growth and market (\SelFirmYears{} firm-years, \SelEvents{} events; every event lies inside the controls' support; Table S3). The $5\times5$ assignment falls back to the size quintile when the book-to-market cell is empty; a strict variant drops those events.

\emph{Survival.} Because labels are censored at fiscal 2025 and disposals arrive at different speeds, we estimate a Cox model of the time from division to first disposal, measured in fiscal years and stratified by filing period $s \in \{2001\text{--}08, 2009\text{--}16, 2017\text{+}\}$,
\begin{equation}\label{eq:cox}
h_i(t) = h_{0,s(i)}(t)\, \exp\!\big(\beta_S S_i + \beta_B\, \mathrm{Block}_i + \boldsymbol{\gamma}'\mathbf{x}_i\big),
\end{equation}
with ties handled by Efron's method and parent-clustered errors; a subsidiary retained through its latest report is censored there. The hazard ratio per standard deviation of $S$ is $\exp(\beta_S\,\sigma_S)$. Two controls accompany the estimate: a permuted signal, whose ratio should cover one, and a synthetic sample with a known ratio that the estimator should recover.

\emph{Reform channels.} The information channel is a difference-in-differences on division announcements by listed parents,
\begin{equation}\label{eq:did}
\mathrm{CAR}_i = \alpha_{t(i)} + \delta_1\, \mathrm{Captive}_i + \delta_2\, \mathrm{Post}_i + \delta_3\, \mathrm{Captive}_i \times \mathrm{Post}_i + \boldsymbol{\gamma}'\mathbf{x}_i + \varepsilon_i ,
\end{equation}
with equity spin-offs as the control, $\mathrm{Post}_i$ equal to one for filings received on or after 1 October 2022, year effects, log market capitalisation and a KOSDAQ indicator, HC1 and parent-clustered errors, a placebo cut at October 2019 inside the pre-period, and a 21-day variant; $\delta_2$ enters on its own because calendar-year effects do not absorb a cutoff inside 2022. The compensation channel regresses the CAR of post-reform \captive{}s on the moneyness $m_i$ of \eqref{eq:moneyness} read from the appraisal price stated in the filing (\AppraisalPrices{} filings; the stated price matches the decree formula within \AppraisalProxyDiff{} percent at the median where the board date lies within a month of the filing), with the formula ratio as a placebo where no right existed, and compares the ex-post payoff of dissenting, $K_i / P_{i,-1} - 1$, with the return to holding the parent for a year. The deterrence and bonding channel classifies the mandatory protection block of every post-reform filing by its content and by the strength of its strongest commitment, and relates both to $S$ and to labelled disposals.

Every analysis was run from a memo that stated the expected range of each estimate before the first run; an estimate outside its range was first checked for code and data errors, and the memos, the log of what was changed and every estimate that moved are in the replication archive. Those that stayed outside their range after checking are reported as they came: the KOSDAQ coefficient of the selection model has the opposite sign to \citeA{kang2025}'s and is insignificant (\SelKosdaq{}, $p$ = \SelKosdaqP{}; Table S3); the signal's AUC spreads more across subgroups than the 0.15 we had allowed (\ValAUCSizeSmall{} for the smallest parents; Table S1); and filings with a charter clause carry a mean signal of \ProtMeanSCharter{}, not the value below 0.40 we expected (Table~\ref{tab:reform}).

\begin{table}[htbp]
\centering
\caption{Research design: alternative explanations and devices}\label{tab:design}
{\scriptsize\singlespacing\setlength{\tabcolsep}{4pt}
\begin{tabular}{>{\raggedright\arraybackslash}p{3.0cm}>{\raggedright\arraybackslash}p{5.4cm}>{\raggedright\arraybackslash}p{6.3cm}}
\toprule
Concern & Alternative explanation & Device (exhibit) \\
\midrule
Model recall & The model remembers the firm's later history & Name- and date-anonymised re-reading; self-reported outside knowledge; subgroup $\Delta$AUC (Sec.~\ref{sec:validity}; Table S1) \\
Measurement bias & The signal is inflated for some kinds of firms & Subgroup AUC and calibration; residual-bias regressions; two-model agreement; external hand-coded benchmark (Table S1) \\
Contribution over hand coding & The signal restates purpose categories or keywords & Nested models M0--M3 with paired-bootstrap increments; keyword and facts-only ladders (Table~\ref{tab:gate1}) \\
Label censoring & Young divisions look ``retained'' & Cox model with period strata, eq.~\eqref{eq:cox}; Kaplan--Meier by tercile (Table~\ref{tab:gate1}, Fig.~S1) \\
Omitted firm condition & A troubled parent writes a disposal-like filing and underperforms & Size, value, momentum, leverage, foreign ownership, board; realised early disposal beside the signal (Tables~\ref{tab:returns}, \ref{tab:block}) \\
Industry; repeat filers & Industry cycles; the same parent files repeatedly & Year + GICS-sector fixed effects; parent- and two-way-clustered errors (Tables~\ref{tab:returns}, \ref{tab:block}, S4) \\
Benchmark & Market-adjusted returns load on size and value & $5\times5$ and $3\times3$ size/B-M portfolios, matched firm, propensity-matched control (Table~\ref{tab:returns}, Panel C) \\
Selection into division & Dividing firms differ from the universe & Firm-year selection logit on the Compustat Global universe; propensity support; matched controls (Table S3) \\
Reform confounded with time & 2023--2026 markets differ & Equity spin-offs as untreated control, eq.~\eqref{eq:did}; placebo cut; 2001--2026, winsorised, median and 21-day variants (Table~\ref{tab:reform}) \\
Reform is a package & The three measures arrived within one quarter & The DiD is read as the package's effect; each channel is tested by the distinct trace it would leave (Table~\ref{tab:reform}, Panels A--C) \\
Right not priced; statements cheap & Appraisal right and protection statements are formalities & Moneyness slope with two placebos; protection-block content vs signal and disposals (Table~\ref{tab:reform}) \\
Label error & Model-read labels are wrong & Tracker cross-check; definition-sensitive labels flipped and dropped (Table S2) \\
\bottomrule
\end{tabular}}
\end{table}
